\documentclass{beamer}
\usepackage[latin9]{inputenc}
\usepackage[french]{babel}
\usepackage{color} 
\usepackage{graphicx} 
\usepackage{listings}
\usepackage{amsmath}
\usepackage{amssymb}
\usepackage{mathptmx}

\definecolor{colKeys}{rgb}{0,0,1} 
\definecolor{colIdentifier}{rgb}{0,0,0} 
\definecolor{colComments}{rgb}{0,0.5,1} 
\definecolor{colString}{rgb}{0.6,0.1,0.1} 

\lstset{%configuration de listings 
float=hbp,% 
basicstyle=\ttfamily\small, % 
identifierstyle=\color{colIdentifier}, % 
keywordstyle=\color{colKeys}, % 
stringstyle=\color{colString}, % 
commentstyle=\color{colComments}, % 
columns=flexible, % 
tabsize=2, % 
frame=trBL, % 
frameround=tttt, % 
extendedchars=true, % 
showspaces=false, % 
showstringspaces=false, % 
numbers=left, % 
numberstyle=\tiny, % 
breaklines=true, % 
breakautoindent=true, % 
captionpos=b,% 
xrightmargin=+0cm, % 
xleftmargin=+0cm
} 


\usetheme{Copenhagen}
\title[Introduction au C avr]{Programmation en C\\Pour avr 8 bits}
\author{J\'er\'emy Cheynet -- INTech \\ Yann Sionneau -- MiNET}
\institute{www.club-intech.fr \\ clubcode.minet.net \\ intlab.minet.net \\ github.com/leroilion/avr}
\date{}
\logo{\includegraphics[height=0.5cm]{image/intech.png}}


\begin{document}
 
\begin{frame}
  \titlepage
\end{frame}

\begin{frame}
  \tableofcontents[]
\end{frame}

\section{Les outils de programmation}
\subsection{Quels sont les outils dont je dispose ?}

\begin{frame}
  \tableofcontents[currentsection]
\end{frame}

\begin{frame}
  \begin{block}{Les diff\'erents outils}
    \begin{itemize}
      \pause
      \item avr-gcc (pour la compilation)
      \pause
      \item avr-objcopy (pour cr\'eer le fichier hex)
      \pause
      \item avrdude (pour flasher)
    \end{itemize}
  \end{block}
\end{frame}

\subsection{Exemple d'utilisation}

\begin{frame}
  \begin{exampleblock}{avr-gcc}
    \pause
    avr-gcc -Wall -mmcu=atmega328p -DF\_CPU=16000000 -c monfichier1.c\\
    \pause
    avr-gcc -Wall -mmcu=atmega328p -DF\_CPU=16000000 -c monfichier2.c\\
    \pause
    avr-gcc -Wall -mmcu=atmega328p -DF\_CPU=16000000 -o monappli monfichier1.o monfichier2.o
  \end{exampleblock}
  \pause
  \begin{exampleblock}{avr-objcopy}
    \pause
    avr-objcopy -O ihex -R .eeprom monappli main.hex
  \end{exampleblock}
  \pause
  \begin{exampleblock}{avrdude}
    \pause
    sudo avrdude -P /dev/ttyUSB0 -c stk500v1 -p m328p -b 57600 -D -U flash:w:main.hex
  \end{exampleblock}
\end{frame}

\subsection{A vous de jouer}

\begin{frame}
  \begin{block}{T\'el\'echarger les sources}
    \pause
    \hyperlink{http://github.com/leroilion/avr}{http://github.com/leroilion/avr}\\
    Et t\'el\'echarger dans les exemples le code blink.c
  \end{block}
  \pause
  \begin{block}{Compiler les sources}
    \pause
    avr-gcc -Wall -mmcu=atmega328p -O2 -DF\_CPU=16000000 blink.c -o blink.out
  \end{block}
  \pause  
  \begin{block}{Cr\'eer le fichier hexad\'ecimal}
    \pause
    avr-objcopy -O ihex -R .eeprom blink.out blink.hex
  \end{block}
  \pause
  \begin{block}{Flasher l'arduino}
    \pause
    sudo avrdude -P /dev/ttyUSB0 -c stk500v1 -p m328p -b 57600 -D -U flash:w:blink.hex\\
    sudo avrdude -c usbtiny -p m328p -U flash:w:blink.hex
  \end{block}
\end{frame}

\section{Hello world !}

\begin{frame}
  \tableofcontents[currentsection]
\end{frame}

\subsection{Structure d'un port}

\begin{frame}
  \pause
  \begin{block}{3 registres}
    \begin{itemize}
      \pause
      \item Le registre DDRx\\
      Registre de configuration du port
      \pause
      \item Le registre PORTx\\
      Registre de sortie du port
      \pause
      \item Le registre PINx\\
      Registre de lecture du port
    \end{itemize}
  \end{block}
\end{frame}

\subsection{Ecrire un 1 ou un 0}

\begin{frame}
  \pause
  \begin{block}{Ecrire un 1 logique}
    \pause
    monport $|$= ( 1 $<$$<$ monbit );\\
    \pause
    \#ifndef sbi\\
    \#define sbi(port,bit) (port) $|$= (1 $<$$<$ (bit))\\
    \#endif
  \end{block}
  \pause
  \begin{block}{Ecrire un 0 logique}
    \pause
    monport \&= $\sim$( 1 $<$$<$ monbit );\\
    \pause
    \#ifndef cbi\\
    \#define cbi(port,bit) (port) \&= $\sim$(1 $<$$<$ (bit))\\
    \#endif\\
  \end{block}
\end{frame}


\subsection{Je te parle}

\begin{frame}
  \pause
  \begin{block}{Configuration en sortie du port}
    \pause
    DDRB $|$= ( 1 $<$$<$ PORTB5 );
  \end{block}
  \pause
  \begin{block}{Ecriture sur un port}
    \pause
    PORTB $|$= ( 1 $<$$<$ PORTB5 ); //Pour mettre le bit 5 du port B \`a 1\\
    \pause
    PORTB \&= $\sim$( 1 $<$$<$ PORTB5 ); //Pour mettre \`a 0 
  \end{block}
\end{frame}

\subsection{Tu me parles}

\begin{frame}
  \pause
  \begin{block}{Configuration du port en entr\'ee}
    \pause
    DDRx \&= $\sim$( 1 $<$$<$ PORTxN );
  \end{block}
  \pause
  \begin{block}{Lecture du port}
    \pause
    PORTx $|$= ( 1 $<$$<$ PORTxN ); //Pour activer le pull--up\\
    PORTx \&= $\sim$( 1 $<$$<$ PORTxN ); //Pour d\'esactiver le pull--up\\
    \pause
    uint8\_t etat = ( PINx \& ( 1 $<$$<$ PINxN ));
  \end{block}
\end{frame}

\subsection{A vous de jouer}

\begin{frame}
  \underline{Objectif:}\\
  Faire un programme qui \'eclaire une LED si un bouton est pouss\'e.\\
  \pause
  \underline{D\'etails techniques:}\\
  Utiliser le PORTB5 en sortie (PORTB7 pour arduino mega), et le PORTB0 en entr\'ee.
  \pause
  \begin{alertblock}{Attention}
    Ne pas oublier le int main() dans le fichier principal.\\
    Penser \`a rajouter l'include standard io.h
  \end{alertblock}
\end{frame}

\lstset{language=c++} 
\lstset{commentstyle=\textit} 
\begin{lstlisting}
#include <avr/io.h>
int main( void )
{
        DDRB |= ( 1 << PORTB5 );
        DDRB &= ~( 1 << PORTB0 );
        PORTB |= ( 1 << PORTB0 );
        while(42)
        {
                if( (PINB & ( 1 << PORTB0 )) )
                        PORTB |= ( 1 << PORTB5 );
                else
                        PORTB &= ~( 1 << PORTB5 );
        }
        return 0;
}
\end{lstlisting} 

\section{C'est bien, mais comment je peux faire un VRAI programme ?}

\begin{frame}
  \tableofcontents[currentsection]
\end{frame}

\subsection{Qu'est-ce qu'un registe ?}

\begin{frame}
  \begin{block}{}
    C'est un octet en m\'emoire\\
    \pause
    qui permet de configurer le microcontr\^oleur.
    \pause
  \end{block}
  \pause
  \begin{exampleblock}{SREG -- AVR Status Register}
    \begin{itemize}
      \pause
      \item I (bit 7) -- Global interrupt enable
      \pause
      \item T (bit 6) -- Copy storage
      \pause
      \item H (bit 5) -- Half carry
      \pause
      \item S (bit 4) -- Sign bit
      \pause
      \item V (bit 3) -- Overflow bit
      \pause
      \item N (bit 2) -- Negative bit
      \pause
      \item Z (bit 1) -- Zero bit
      \pause
      \item C (bit 0) -- Carry
    \end{itemize}
  \end{exampleblock}
\end{frame}

\subsection{Le fil rouge}

\begin{frame}
  \underline{Objectif:}\\
  Faire un programme qui fait clignoter une led en utilisant le TIMER1 sur 16 bits.\\Pour cela, on fera:\\
  \begin{itemize}
    \pause
    \item On activera les interruptions d'overflow du TIMER1
    \pause
    \item On fera compter le TIMER1 pour avoir une interruption toutes les secondes.
    \pause
    \item On regardera l'\'etat de la PIN associ\'ee \`a la LED pour le changer.
  \end{itemize}
\end{frame}


\subsection{Le datasheet: la bible du programmateur bas niveau}

\begin{frame}
  \hyperlink{http://github.com/leroilion/avr/tree/master/datasheet/Atmegaxx8.pdf}{http://github.com/leroilion/avr}
  \pause
  \begin{block}{TCCR1A -- TCCR1B (page 132 -- 134)}
    Mode normal (WGMx = 0 ), Source d'horloge (CSx = 101), Pas de comparaison (COMx = 0) 
  \end{block}
  \pause
  \begin{block}{TCNT1H -- TCNT1L (page 136) }
    Incr\'ementation toute les 64$\mu$s \pause $\Longrightarrow$ Compter jusqu'\`a 15625 \pause $\Longrightarrow$ Mettre 49910 = 65535 - 15625 dans TCNT.
  \end{block}
  \pause
  \begin{block}{TIMSK1 (page 137)}
    Activer l'interruption d'overflow \pause $\Longrightarrow$ Activer TOIE1
  \end{block}
\end{frame}

\subsection{Tu me vois, tu me vois plus}

\lstset{language=c++} 
\lstset{commentstyle=\textit} 
\begin{lstlisting}
#include <avr/io.h>
#include <avr/interrupt.h>
...
TCCR1A = 0bxxxxxx00;
TCCR1B = 0bxxx00101;
TIMSK1 |= ( 1 << TOIE1 );
...
ISR( TIMER1_OVF_vect)
{ ...
\end{lstlisting}

\subsection{A vous de jouer}

\begin{frame}
  \underline{Objectif:}\\
  Faire un programme qui permet de contr\^oler la luminosit\'e d'une LED en utilisant le PWM.\\
  \pause
  \underline{D\'etails techniques:}\\
  \begin{itemize}
    \pause
    \item Utiliser le port B5 pour la sortie de la LED (PORTB7 pour arduino mega).
    \pause
    \item Utiliser le TIMER de votre choix en mode PWM
  \end{itemize}

  \pause
  \begin{alertblock}{Attention}
    Ne pas oublier l'include $<$avr/interrupt.h$>$\\
    Penser \`a g\'erer \emph{TOUS} les vecteurs d'interruption
  \end{alertblock}
\end{frame}

\section{Les pi\`eges \`a \'eviter}

\begin{frame}
  \begin{alertblock}{Les pi\`eges}
    \pause
    \begin{itemize}
     \item \'Economiser la m\'emoire (probl\`eme de la pile)
     \pause 
     \item Faire attention avec les \emph{float}, les . et les \emph{double}
     \pause
     \item \'Economiser la puissance de calcul (calcul en 8 bits)
     \pause
     \item Faire attention \`a l'overflow
     \pause
     \item Rajouter l'option \emph{volatile} devant les variables
    \end{itemize}

  \end{alertblock}
\end{frame}


\section*{Conclusion}

\begin{frame}
  Nous avons vu :\\
  \begin{itemize}
    \pause
    \item Les outils de programmation
    \pause
    \item Comment configurer, lire et \'ecrire sur un port
    \pause
    \item Ce qu'est un registre, et comment le configurer \`a l'aide du datasheet
    \pause
    \item L'utilisation des interruptions
  \end{itemize}
  \pause 
  Des exemples simple :\\
   \begin{itemize}
    \pause
    \item \hyperlink{http://github.com/leroilion/avr}{http://github.com/leroilion/avr}
   \end{itemize}
  Bibliographie :\\
    \begin{itemize}
      \pause
      \item \hyperlink{http://www.tavernier-c.com/livres.htm}{\underline{Microcontr\^oleurs AVR : des ATtiny aux ATmega} de Christian Tavernier }
    \end{itemize}

\end{frame}

\end{document}
